\chapter{Resolución de problemas}\label{ch:g2:problems}

Una historia con números esconde una pregunta: ¿sumar? ¿llevar? ¿multiplicar?
Este capítulo ofrece una manera de decidir, utiliza imágenes de barras simples para
\emph{ver} la operación y lee información de tablas y
pictogramas.

\section{Eligiendo la operación}

\begin{method}[¿Qué operación?]\label{met:g2:problems:choose}
\begin{enumerate}
  \item \emph{juntando, en todo}: agregar ($+$);
  \item \emph{quitando, cuantos quedan, cuantos mas}: restar
        ($-$);
  \item \emph{tantos grupos iguales}: multiplicar ($\times $);
  \item Haz un dibujo rápido cuando no estés seguro y responde siempre con un
        frase.
\end{enumerate}
\end{method}

\begin{omfigure}
\begin{tikzpicture}[scale=0.6]
  % put together
  \draw[thick, fill=omDef!20] (0,0) rectangle (3,0.8);
  \draw[thick, fill=omThm!20] (3,0) rectangle (4.6,0.8);
  \node[font=\small] at (1.5,0.4) {$24$};
  \node[font=\small] at (3.8,0.4) {$15$};
  \draw[Stealth-Stealth] (0,-0.4) -- (4.6,-0.4);
  \node[below, font=\small] at (2.3,-0.5) {? in all: $24 + 15$};
  % compare
  \begin{scope}[xshift=8cm]
  \draw[thick, fill=omDef!20] (0,0.9) rectangle (4.6,1.7);
  \node[font=\small] at (2.3,1.3) {$40$};
  \draw[thick, fill=omThm!20] (0,0) rectangle (2.8,0.8);
  \node[font=\small] at (1.4,0.4) {$28$};
  \draw[Stealth-Stealth] (2.8,0.4) -- (4.6,0.4);
  \node[below, font=\small] at (2.3,-0.2) {? more: $40 - 28$};
  \end{scope}
\end{tikzpicture}

{\small Imágenes de barras: dos barras unidas significan una \omterm{def:g1:addition:def}{suma}; dos barras
comparado significa una resta.}
\end{omfigure}

\begin{example}[Un problema de dos pasos]\label{ex:g2:problems:twostep}
``Sam compra $3$ paquetes de $6$ pegatinas, luego da $5$ lejos. cuantos
¿Se mantiene? ''
\begin{enumerate}
  \item Grupos: $3 \times 6 = 18$ pegatinas;
  \item donar: $18 - 5 = 13$.
\end{enumerate}
Sam mantiene $13$ pegatinas. Vuelva a leer la pregunta: ¿la historia es realmente?
terminado? Aquí si.
\end{example}

\begin{example}[Un número extra oculto]\label{ex:g2:problems:extra}
``Un $210$-libro de páginas\,\dots'' --- cuidado, no todos los números en un
Se necesita historia. ``Ben tiene $8$ rojo y $6$ canicas azules y es $7$
años. ¿Cuántas canicas?'' La edad es un señuelo: $8 + 6 = 14$
canicas. Ordena los números antes de calcular.
\end{example}

\section{Tablas y pictogramas}

\begin{example}[una mesa]\label{ex:g2:problems:table}
Pasteles vendidos en un puesto:
\begin{center}
\renewcommand{\arraystretch}{1.3}
\begin{tabular}{l|ccc}
 & Saturday & Sunday & total \\
\hline
muffins & $12$ & $9$ & $21$ \\
tarts & $7$ & $10$ & $17$ \\
\end{tabular}
\end{center}
Lectura: tartas dominicales, $10$. Informática: magdalenas durante los dos días,
$12 + 9 = 21$; más tartas el domingo que el sábado $10 - 7 = 3$.
\end{example}

\begin{example}[un pictograma]\label{ex:g2:problems:pictogram}
Lea primero la clave: aquí una manzana significa $2$ trozos de fruta.

\begin{omfigure}
\begin{tikzpicture}[scale=0.62]
  \node[left, font=\small] at (0,2) {red};
  \foreach \i in {0,1,2,3} \fill[omThm] (0.5+\i*0.75,2) circle (0.24);
  \node[left, font=\small] at (0,1) {green};
  \foreach \i in {0,1} \fill[omMeth] (0.5+\i*0.75,1) circle (0.24);
  \node[left, font=\small] at (0,0) {yellow};
  \foreach \i in {0,1,2} \fill[omProp] (0.5+\i*0.75,0) circle (0.24);
  \node[right, font=\small] at (4.0,1) {key: one dot $= 2$ fruits};
\end{tikzpicture}

{\small Un pictograma: el rojo muestra $4$ puntos, entonces $4 \times 2 = 8$ rojo
frutas; verde $2$ puntos, $4$ frutas; amarillo $3$ puntos, $6$ frutas.}
\end{omfigure}
\end{example}

\section{Ejercicios}

\begin{exercise}[$\star $]\label{exo:g2:problems:1}
Nombre la operación (no calcule): ``$34$ y $28$ canicas puestas
juntos''; ``$50$ dulces, $12$ comido''; ``$4$ cajas de $5$
manzanas''; ``Ana tiene $30$, Leo tiene $18$: ¿Cuantos más tiene Ana?''.
\end{exercise}

\begin{exercise}[$\star $]\label{exo:g2:problems:2}
Resuelva: ``Un aparcamiento tiene $156$ coches; $48$ dejar. ¿Cuántos se quedan?
\end{exercise}

\begin{exercise}[$\star $]\label{exo:g2:problems:3}
Resuelva: ``$5$ niños cada elección $4$ flores. cuantas flores hay
¿Todos?
\end{exercise}

\begin{exercise}[$\star $]\label{exo:g2:problems:4}
Haz un dibujo de una barra y luego resuelve: ``Tom tiene $45$ tarjetas, Mia tiene $17$
menos. ¿Cuántos tiene Mía?
\end{exercise}

\begin{exercise}[$\star $]\label{exo:g2:problems:5}
Dos pasos: ``Una caja contiene $6$ botellas. una tienda consigue $4$ cajas y
vende $9$ botellas. ¿Cuántas botellas quedan?
\end{exercise}

\begin{exercise}[$\star $]\label{exo:g2:problems:6}
Encuentra y tacha el número inútil, luego resuelve: ``Ben es $8$
años y tiene $23$ canicas. el gana $15$ más. cuantas canicas
¿ahora?''
\end{exercise}

\begin{exercise}[$\star $]\label{exo:g2:problems:7}
Usando la tabla de \cref{ex:g2:problems:table}: cuantas tartas se vendieron
durante los dos días? ¿Cuántas tartas de todo tipo el sábado?
\end{exercise}

\begin{exercise}[$\star $]\label{exo:g2:problems:8}
Usando el pictograma de \cref{ex:g2:problems:pictogram}: cuantos
frutas en total? ¿Qué color tiene más?
\end{exercise}

\begin{exercise}[$\star $]\label{exo:g2:problems:9}
Resuelva: ``Una cinta es $80$ cm de largo. Lea corta $35$ centímetro. cuanto dura
¿La pieza que queda?
\end{exercise}

\begin{exercise}[$\star\star $]\label{exo:g2:problems:10}
``Un panadero hace $5$ bandejas de $4$ pasteles por la mañana y $3$ bandejas de $4$ pasteles por la tarde. ¿Cuántas tortas ese día?'' (Dos
multiplicaciones, entonces un \omterm{def:g1:addition:def}{suma} --- ¡o contar las bandejas primero!)
\end{exercise}
